\newcommand{\MainEvidence}{Proposed LIOM accepts 29--30 of 30 scenes per repetition, with 29 accepted in all repetitions; mean wall time is 7.83 s. Hard RSC--iterative accepts 28--30 per repetition and averages 4.03 s.}
\newcommand{\PhysicalEvidence}{For crossing, proposed LIOM accepts 10/10, 5/10, 1/10, versus 10/10, 6/10, 1/10 for iterative hard RSC, across increasing difficulty levels. For passage, proposed LIOM accepts 0/5, 0/5, 0/5, versus 0/5, 0/5, 0/5 for iterative hard RSC, across increasing difficulty levels. For obstacles, proposed LIOM accepts 5/5, 2/5, 4/5, versus 4/5, 4/5, 3/5 for iterative hard RSC, across increasing difficulty levels.}
\newcommand{\RecoveryEvidence}{The targeted passage study reports proposed / iterative hard-RSC acceptance as follows: $N=5$, width 2.6 m: 1/3 / 0/3; $N=5$, width 2.0 m: 0/3 / 0/3; $N=8$, width 2.6 m: 0/3 / 0/3; $N=8$, width 2.0 m: 0/3 / 0/3; $N=10$, width 2.6 m: 0/3 / 0/3; $N=10$, width 2.0 m: 0/3 / 0/3. All 18 geometries and both outcomes are retained.}
\newcommand{\SensitivityEvidence}{The one-factor study retains all six outcomes at every setting. The table reports mean times; its caption states that all six anchor tasks are accepted at each tested setting. The experiment uses one-factor changes.}
\newcommand{\TrackingEvidence}{Nominal tracking evaluates 35 motions; maximum error is 0.132 m, maximum terminal error is 0.105 m, and 0 runs record contact. Perturbed tracking evaluates 35 motions; maximum error is 0.139 m, maximum terminal error is 0.121 m, and 0 runs record contact.}
\newcommand{\HardIterationEvidence}{On the main benchmark, iterative hard boxes accept 18 scenes per repetition, compared with 16 for one-shot hard boxes.}
\newcommand{\HardEntryEvidence}{Soft LIOM averages 8.42 s and 2.29 QP calls; the proposed entry policy averages 7.83 s and 2.46 calls.}
\newcommand{\AcceptanceEvidence}{The six anchor tasks are accepted under all four recovery-criterion/feedback variants. Across the eight frozen suites, 23 reconstructed soft candidates fail a physical check, in 19 attempts; 15 of those attempts later return an accepted motion within budget. The exact-finishing branch returns 9 accepted motions.}
\newcommand{\EnvelopeEvidence}{The envelope table reports computation on the same six tasks, all accepted under the tested envelopes.}
\newcommand{\FailureEvidence}{Physical-task terminal outcomes are adaptive: accepted 27, accepted over budget 2, budget 30, initialization 1; liom: accepted 30, budget 28, accepted over budget 1, initialization 1; hard--iter: accepted 12, budget 47, initialization 1; rsc--iter: accepted 28, budget 26, accepted over budget 5, initialization 1.}
\newcommand{\SolverEvidence}{The OSQP numerical audit uses identical matrices from three scenes for soft boxes, hard boxes, and hard RSC. Readiness stopping and ordinary termination at two tolerances are reported separately.}
\newcommand{\InitializationEvidence}{Across the 1,212 attempts in the eight frozen suites, four initialization failures concern one physical scene evaluated by four methods. The static-guide fallback follows 93 SFM simulation timeouts; no scalar-timing failure is recorded in these suites. Counts refer to attempted executions, including timing repetitions.}
\section*{Reviewer 1}
\reviewitem{R1.1 --- Team-size scalability}{Discuss computational scalability beyond seven UAVs.}{We added physical-congestion tasks with 10, 15, and 20 UAVs and a separate study of $N\in\{5,10,15,20\}$ and $K\in\{30,60,90\}$. Nominal corridor-growth limits are retained. A box-distance test removes a pair restriction only when the spatial boxes already ensure the protected separation; all pairs remain in final verification. Section 6.4 separates assembly, screening, QP solution, and checking costs. Section 7.6 reports measured scaling and stage times. \PhysicalEvidence}
\reviewitem{R1.2 --- Advantage over hard RSC}{Clarify when recovery helps, given the faster hard RSC result.}{The revised method first tries a bounded hard separating-plane QP and invokes the relaxed loop if that attempt does not yield an accepted motion. We compare one-shot and iterative hard formulations. The new winding-counterflow study identifies an independently verified conflict between initial relative restrictions and exact kinematics. In the initial three paired repetitions, the representative task is recovered three times while iterative hard RSC accepts none; a separate 12-realization check accepts seven versus six tasks. A physical-width/team/horizon map locates these conditions with unchanged corridor-growth limits. The paired external-baseline comparison returns accepted Winding A motions with LIOM and hard RSC in 26.36 and 25.14 s. The nonlinear comparisons were rerun with Ipopt/ma57: ASCO ends without an accepted Winding A motion at 33.72 s, while it accepts C27 in 8.61 s. The case table reports these outcomes individually. Winding B returns a LIOM motion at 31.92 s, explicitly marked as late. Both earlier and later records are retained; the response distinguishes fixed-corridor incompatibility from outer-loop outcomes. \MainEvidence\ \RecoveryEvidence}
\reviewitem{R1.3 --- Parameter sensitivity}{Explain and test tolerances, penalty weights, and expansion limits.}{Section 7.2 explains the penalty scale, conditioning cap, and reference scale for normalized recovery. Section 7.7 varies one factor at a time on six fixed tasks: effort weight, initial penalties, caps, growth factor, defect/slack tolerances, corridor limits, candidate-readiness tolerance, and principal SFM parameters. Physical comparisons retain nominal corridor limits. \SensitivityEvidence}
\reviewitem{R1.4 --- Length and consistency}{Tighten the review and methods, and check terminology.}{Related Work now focuses on corridor optimization, local coordination, and the fixed-horizon task. Repeated novelty claims were removed. Variables are defined at first use, and terminology is consistent across the flowchart, pseudocode, figures, and tables. The algorithm explicitly shows the hard entry, recovery loop, acceptance feedback, and failure exits.}
\reviewitem{R1.5 --- Control interface}{Discuss tracking and the suggested sine-resistance-network study.}{Section 7.8 describes a shared execution clock and individual tracking controllers. New MuJoCo simulations include rigid-body attitude dynamics, bounded rotor forces, motor lag, drag, and disturbances. Huang et al.'s suggested study is cited for its planning/control simulation context, with its vehicle-specific mechanism distinguished from the UAV formulation. \TrackingEvidence}
\clearpage
\section*{Reviewer 2}
\reviewitem{R2.1 --- Completion time}{Define $T_f$.}{Section 3 defines $T_f$ as the prescribed common physical flight-completion time, with $h=T_f/K$. It is distinct from the computation budget. Fixed endpoint states apply at zero and $T_f$.}
\reviewitem{R2.2 --- Cartesian limits}{Explain identical limits in three coordinates.}{The generic hovering-capable multirotor uses isotropic Euclidean translational planning limits. This assumption does not identify equal physical actuator authority in every direction. The old cube is replaced by a richer polyhedral inner approximation, containing the cube and reaching the full axial Euclidean limits while retaining linear constraints. Final checks enforce Euclidean bounds. Sections 3.2 and 7.7 define and compare the motion envelopes.}
\reviewitem{R2.3 --- Collision notation}{Explain $c$, $t$, and $\mathbf p_i(t)$.}{Section 3 defines surface clearance $c$, footprint radius $r$, physical time $t\in[0,T_f]$, and UAV-center position $\mathbf p_i(t)$. Pair distance must exceed $2r+c$ and center-to-obstacle distance must exceed $r+c$. Workspace and obstacle-distance notation are defined alongside these conditions.}
\reviewitem{R2.4 --- Objective}{Explain $w_a$ and $C(V)$.}{Section 3.4 defines $C(V)=w_a h\sum_{i,k}\|\mathbf a_{i,k}\|^2$ as weighted integrated squared translational acceleration. Acceleration is obtained from consecutive velocities and is constant within an interval. The weight regularizes the motion; the objective is not described as electrical energy. The comparison centers on computation and independently checked recovery cases; complete task-level outcomes remain available. Units and weight sensitivity are stated explicitly.}
\reviewitem{R2.5 --- Reachable speed}{Explain $b_{i,k}$ and simplify zero endpoint velocities.}{The formula is simplified to $b_{i,k}=\min\{v_{\max},a_{\max}t_k,a_{\max}(T_f-t_k)\}$. It bounds speed reachable from rest and compatible with braking to rest at the fixed final time. This bound supplies the inter-node displacement padding in corridor construction.}
\reviewitem{R2.6 --- Complexity}{Provide complexity.}{Section 6.4 gives $O(NK+M)$ sparse assembly/storage for $M$ active pair--node triples, $O(N^2K)$ worst-case screening, $O(NKB)$ spatial construction under fixed growth rules for $B$ obstacles, and $O((N^2+NB)K)$ complete-interval checking. QP solution also depends on sparse factorization, conditioning, and iterations. The $N$--$K$ study reports measured dimensions and stage times.}
\reviewitem{R2.7 --- Practical cooperation}{Explain how coordination is implemented.}{All assigned starts/goals and the shared static map enter a centralized solve. Relative restrictions couple UAV positions at synchronized nodes. The accepted joint trajectories are distributed to individual tracking controllers with a shared start clock. Section 7.8 describes and simulates this separation between centralized coordination and local thrust/attitude feedback.}
\reviewitem{R2.8 --- UAV type and parameters}{Specify the UAV type and simulation parameters.}{The paper explicitly studies hovering-capable rotary-wing UAVs. Section 7.2 gives planning radius, clearance, motion limits, horizon, and grid. Section 7.8 gives generic quadrotor mass, inertia, rotor geometry, force limits, motor lag, integration step, and disturbances. These define a reproducible simulated aircraft rather than identifying a commercial platform.}
\section*{Reviewer 3 --- Major Comments}
\reviewitem{R3.M1 --- Relation to submitted references [20] and [31]}{Identify inherited components and additional contributions; include the closest task-matched baseline.}{We reread both originals and revised the attribution. Iterative corridor construction, external penalties, and separate kinematic/separation penalties are acknowledged as antecedent ideas. The contribution list now focuses on the fixed common-time rest-to-rest formulation, affine QP structure, reachable-speed padding, endpoint moment correction with exact rollout, and the new hard-entry and acceptance-feedback procedure. The nonemptiness proof explains this local formulation; it is not presented as a general convergence theorem. Table 1 and the final Related Work paragraph give these distinctions. Hard box--iterative and Soft LIOM use the same UAV task, reference, corridors, reconstruction opportunities, and acceptance rules, providing controlled task-matched adaptations of the closest mechanisms. The revised comparison also includes task-matched ASCO, direct nonlinear programming, and sequential convex programming, with common endpoint, motion, initialization, and final-acceptance conditions. The new manuscript subsection ``Task-Matched Nonlinear and Sequential Convex Comparisons'' presents the paired case results. The ASCO risk-range update follows the original paper and archived implementation. SCP uses the original distance-norm tangent with sparse numeric OSQP assembly. The ten fixed tasks include ordinary crossings/passages, two winding layouts, and 15- and 20-UAV tasks. The table reports each task, its computation time, physical acceptance, and any budget overrun; no population success-rate ranking is inferred.}
\reviewitem{R3.M2 --- Physical difficulty}{Add denser teams, narrower passages, and multiple random instances.}{The new 60-task physical benchmark includes ten randomized crossing layouts at each of 10, 15, and 20 UAVs, plus five layouts at each of three passage/obstacle levels with ten UAVs. The workspace is $16\times16\times7$ m; passage openings are 3.2, 2.6, and 2.0 m. The design and seeds were saved before execution, and all outcomes are retained. These tests change physical congestion while preserving nominal corridor-growth settings. A further winding-counterflow study uses three alternating physical openings. In the initial three paired repetitions, the representative five-UAV task is accepted by the complete method and not by iterative hard RSC. The additional external-baseline study also records a hard-RSC acceptance of that task in 25.14 s. The original positive phase-I certificate applies to its initial corridors, and reconstruction can change their compatibility. A separate set of 12 previously untested realizations yields seven versus six accepted tasks. A physical-width/team/horizon map shows the additional recovery region, with every attempted outcome retained. \PhysicalEvidence}
\reviewitem{R3.M3 --- Equal reconstruction opportunities}{Give hard methods the same outer loop; the growth-cap sweep confounds iteration with softening.}{Hard box--iterative and Hard RSC--iterative now use the same reconstruction opportunities, reference-update rules, and static-guide policy. A usable numerical iterate can guide the next construction even when it is not accepted as a flight; otherwise the shared spatial-center/guide fallback applies. Every accepted motion passes identical checks. One-shot hard variants isolate the iteration restriction. Corridor-size changes are assessed as parameter sensitivity, separately from physical difficulty. In the new winding task, the initial hard RSC phase-I LP has matching positive primal and dual objective 0.21915. Removing only the relative rows or only the kinematic equalities makes that diagnostic LP feasible at zero violation. Thus the example identifies the coupled conflict while giving the hard baseline the same reconstruction opportunities. \HardIterationEvidence}
\reviewitem{R3.M4 --- Hard-first strategy and exact penalties}{Avoid multiple relaxed solves when an initial hard QP suffices; discuss finite exact-penalty thresholds.}{The proposed method begins with a bounded hard RSC-QP attempt. If its reconstructed motion passes checking, it returns immediately; otherwise the relaxed box-corridor loop follows. Soft LIOM removes this entry step. A further controlled study varies the entry allowance between 1, 3, and 5 s on ten tasks with identical references and a 30 s total budget. For C29, the 3 s allowance reduces total planning time from 19.27 to 3.63 s and QP calls from five to one. The 3 s setting is recommended for new runs; archived comparisons retain their explicitly labeled 1 s setting. Table~\ref{tab:entry-budget-response} reproduces the full ablation.  Section 6.2 discusses the linear slack penalty: under the usual convex-QP duality assumptions, a sufficiently large weight relative to optimal multipliers can be exact for those inequalities. A quadratic defect penalty generally remains inexact at finite weight. Caps therefore address conditioning, while measured recovery and final checking determine acceptance. No universal finite exactness threshold is asserted for both terms. \HardEntryEvidence}
\begin{table}[htbp]

\caption{Hard-entry budget ablation. Each cell gives total planning time in seconds and QP calls. All three settings share the recorded reference and a 30 s total budget, with rotated execution order. $\times$ marks no accepted motion and $\dagger$ marks late acceptance.\label{tab:entry-budget-response}}
\small\begin{tabularx}{\textwidth}{Xrrr}\toprule
Task & 1 s entry: time / QPs & \textbf{3 s entry: time / QPs} & 5 s entry: time / QPs\\\midrule
C01 & 1.75 / 1 & 0.95 / 1 & 0.88 / 1\\
C03 & 1.35 / 1 & 1.44 / 1 & 1.46 / 1\\
C05 & 2.84 / 1 & 2.91 / 1 & 2.89 / 1\\
C25 & 1.20 / 1 & 1.18 / 1 & 1.17 / 1\\
C27 & 24.08 / 6 & 26.15 / 6 & 28.29 / 6\\
C29 & 19.27 / 5 & 3.63 / 1 & 3.65 / 1\\
Winding A & 28.70 / 5 & 30.00 / 5 & 31.33$^{\dagger}$ / 5\\
Winding B & 30.48$^{\dagger}$ / 6 & $\times$\,30.07 / 5 & $\times$\,30.04 / 5\\
15-UAV & $\times$\,30.08 / 7 & 29.09 / 6 & $\times$\,30.66 / 6\\
20-UAV & 28.37 / 6 & 27.04 / 6 & 30.35$^{\dagger}$ / 6\\
\bottomrule\end{tabularx}
\end{table}

\reviewitem{R3.M5 --- Reconstructed-motion acceptance}{Clarify the candidate/output gap, add feedback after failed checks, and normalize recovery.}{The final output comes from corrected controls and exact integration. Its independent checker examines every obstacle and UAV pair, endpoints, workspace, and Euclidean speed/acceleration limits. A numerically recovered candidate that fails this check triggers a tighter normalized defect tolerance and a capped defect-weight increase, followed by an exact-kinematics, zero-slack attempt in the current boxes. If that also fails, reconstruction continues within budget. The normalized defect is $\bar D=(h_0/h)^2D/(NK)$; slack is averaged over active pair--node restrictions. These are numerical criteria, not substitutes for physical acceptance. The acceptance-feedback ablation evaluates this combined procedure and records rejected-check reasons. \AcceptanceEvidence}
\reviewitem{R3.M6 --- Motion-envelope conservatism}{Replace or justify the inscribed cube and use physical Euclidean bounds for final checks.}{We adopted a richer polyhedral inner approximation to retain linear QP constraints. The nominal convex hull of six unit-axis directions and eight normalized cube vertices has 24 facets. It contains the old cube, reaches full axial authority, and lies inside the Euclidean ball. Scaling gives velocity and acceleration limits. Final checks now use Euclidean norms. The original cube and a 26-vertex, 48-facet envelope are compared with the nominal envelope. \EnvelopeEvidence}
\reviewitem{R3.M7 --- Theory and termination statistics}{Nonemptiness does not prove convergence; report rounds and failure branches.}{The proposition is explicitly titled ``Nonemptiness of the relaxed fixed-corridor subproblem.'' It assumes nonempty, endpoint-compatible spatial boxes and establishes a local feasible point. Numerical recovery and final physical acceptance are defined separately. No general outer-convergence or bounded-round success claim is made. Tables give outer rounds and QP calls. Records retain initialization/construction failures, numerical nonrecovery, rejected reconstructed motions, stalled hard reconstruction, and budget exits. \FailureEvidence}
\reviewitem{R3.M8 --- QP solver and timing decomposition}{Use a dedicated solver, compare matched QPs, and report stage and scale effects.}{Sparse numeric QPs now use OSQP. End-to-end timing includes reference generation, corridors, assembly, solver setup/data exchange, iterations, endpoint reconstruction, and verification. CPU time includes MATLAB and the worker. A separate identical-matrix OSQP audit compares candidate-readiness stopping with ordinary solver termination at two tolerances. All planning QPs use OSQP. The additional ASCO and direct-NLP comparisons use Ipopt with ma57 for their nonlinear collision constraints. Both baselines are rerun on all ten fixed tasks with the same inputs and budgets. The $N$--$K$ study and stage plot address scale effects. \SolverEvidence}
\reviewitem{R3.M9 --- Physical realism}{Add higher-fidelity tracking or state the model scope.}{MuJoCo simulations include six-degree-of-freedom rigid-body dynamics, four bounded rotor forces, motor lag, gravity, drag, and collision geometry. Position/velocity feedback drives a rotation-matrix attitude controller. Fixed controller settings are used under nominal and perturbed mass, lag, and wind. Terminal error is measured at $T_f$; the logged settling interval is separate. The Conclusion identifies the static-map setting and additional modeling needed for downwash and perception uncertainty. The complete controller and aircraft simulation execute in software; no flight-control hardware is connected. \TrackingEvidence}
\reviewitem{R3.M10 --- Initializer effects}{Assess SFM sensitivity/failure and separate reference geometry from later guide availability.}{The initializer study crosses SFM, static-guide, and straight geometry with guide availability during relocation. The guide-only fallback after SFM/timing failure is explicitly recorded. The sensitivity study varies repulsion, decay length, look-ahead, response time, and tangential force independently. All initialization failures and fallback reasons remain in the saved records. \InitializationEvidence}

\section*{Reviewer 3 --- Minor Comments}
\reviewitem{R3.m1 --- Quantitative abstract}{Add sizes, accepted counts, and timing.}{The abstract states the evaluation scale of up to 20 UAVs and describes the two recovery layers and closed-loop tracking study. Detailed runtimes and accepted-case outcomes are presented in the results tables and in this response. \MainEvidence}
\reviewitem{R3.m2 --- Plain language and terminology}{Reduce dense prose and distinguish meanings of recovery.}{The Introduction, Related Work, and Conclusion have been tightened. ``Numerical recovery'' means meeting normalized defect/slack criteria. Corridor rebuilding and exact control integration are specified explicitly. ``Acceptance'' means passing the complete-interval physical checks. The local nonemptiness result is not described as an overall planning-success guarantee.}
\reviewitem{R3.m3 --- Units and $w_a$}{Explain mixed units and selection of $w_a=0.05$.}{The weights carry the reciprocal units appropriate to their terms: the effort integral has units $\mathrm{m}^2/\mathrm{s}^3$, $D/h$ has units $\mathrm{m}^2/\mathrm{s}$, and the slack sum has units m. The nominal $w_a=0.05$ regularizes motion and is varied by half/double in sensitivity tests; the coefficient weights acceleration effort. One positive scale divides the whole objective before solution, preserving relative weights and minimizers.}
\reviewitem{R3.m4 --- Reference resampling}{Explain the 0.02 s to 0.2 s mapping.}{Section 5.2 now gives the implemented arc-length mapping. SFM samples form polylines; repeated points are removed. Each route receives an admissible scalar rest-to-rest timing law over the common horizon and is evaluated on the optimization grid. The outgoing tangent is used at an interior geometric knot and the incoming tangent at the terminal knot. The prescribed timing law therefore determines the resampling locations along each route. Full vector-dynamic consistency at corners is addressed by optimization.}
\reviewitem{R3.m5 --- Scalar timing failures}{Report failures after path-length scaling.}{Both scalar motion limits are checked after scaling. An SFM timing failure invokes the recorded guide-only fallback; a remaining timing failure terminates initialization. The records distinguish an SFM exception from failure of the complete initialization policy. \InitializationEvidence}
\reviewitem{R3.m6 --- Trajectory-panel size}{Use fewer larger panels and move the complete grid out of the body.}{The main figure uses three representative scenes with matched views of reference and accepted motion, metric axes, and endpoint markers. A projection/control-profile figure adds planar geometry, speed, and acceleration histories. The full 30-scene grid is retained in the reproducibility material.}
\reviewitem{R3.m7 --- Timing-marker overlap}{Use log scale or offsets.}{The timing figure uses a logarithmic axis and horizontal method offsets. Marker shapes distinguish methods, crosses mark scene--method combinations with at least one unsuccessful timed repetition; all attempts enter the runtime summaries. Proposed LIOM is explicitly named and consistently emphasized.}
\reviewitem{R3.m8 --- Grayscale readability}{Add symbols or hatching.}{Method plots use distinct marker shapes, and stage-time bars combine hatching with color. Outcome symbols and captions explain the encodings.}
\reviewitem{R3.m9 --- Outer iterations}{Add outer-round counts to the comparison table.}{The main table gives mean outer rounds and mean QP calls. These differ because the hard entry and exact-kinematics finish can add solves within a round. Every attempt also retains its phase history.}
\reviewitem{R3.m10 --- Reproducibility}{Release code and scenario definitions.}{The code, scenario definitions, simulator/controller, analysis scripts, manuscripts, and complete raw outcomes are publicly available at \url{https://github.com/libai1943/uav_liom}. Large experiment archives, including unsuccessful attempts and the earlier records, are included in the repository release. Source hashes and file manifests identify the tested snapshots. The recommended configuration uses a 3 s hard entry, and the archived 1 s experiments remain explicitly identified.}

